\documentclass[../../../include/open-logic-section]{subfiles}

\begin{document}

\olfileid[ar]{sfr}{infinite}{card-sb}

\olsection{ملحق: إثبات مبرهنة شرودر–برنشتاين}

بقي قبل مفارقة نظرية المجموعات الساذجة برهان أخير، ساذج في مبانيه دقيق في صنعته: برهان مبرهنة شرودر–برنشتاين
(\olref[sfr][siz][sb]{thm:schroder-bernstein}). ومفادها أنه إذا كان
$\cardle{A}{B}$ و$\cardle{B}{A}$، كان $\cardeq{A}{B}$؛ أي إنه متى وُجدت
!!{injection}s $f \colon A \to B$ و$g \colon B \to A$، وُجدت
!!a{bijection} $h \colon A \to B$.

جرى البحث في هذا الفصل على مفهوم ديديكند في \emph{الإغلاقات}. وبهذا المفهوم أقام ديديكند برهانًا حسنًا لمبرهنة شرودر–برنشتاين، نعرضه الآن. وطريقتنا قريبة من معالجة
\citet[صص.~157--8]{Potter2004}؛ فمن شاء أن يرجع إليها وجد عرضًا يختلف قليلًا ويتفق في الأصل. ويكفي قليل من البحث في غوغل ليتبين أن لهذه المبرهنة براهين
\emph{كثيرة} مختلفة طريفة، كما لعدم نسبية
$\sqrt{2}$ براهين من هذا القبيل.

وباستخدام ترميز شبيه بما في
\olref[sfr][infinite][dedekind]{Closure}، لتكن $E$ مجموعة، ولتكن
$B\subseteq E$ و$f\colon E\to E$، ونضع
\[
\Closureofunder{f}{B} = \bigcap \Setabs{X}{X\subseteq E\text{ و}B \subseteq X
\text{ و$X$ منغلقة تحت $f$}}
\]
وبحسب هذا التعريف، تكون
$\Closureofunder{f}{B}$ أصغر مجموعة منغلقة تحت $f$ وتحتوي~$B$، بمعنى
ما يأتي:

\begin{lem}\ollabel{Closureprops}
	لكل مجموعة $E$، ولكل دالة $f\colon E\to E$، ولكل $B\subseteq E$:
	\begin{enumerate}
		\item\ollabel{Closurehaselem} $B \subseteq \Closureofunder{f}{B}$؛ و
		\item\ollabel{Closureclosed} $\Closureofunder{f}{B}$ منغلقة تحت $f$؛ و
		\item\ollabel{Closuresmallest} إذا كانت $X$ منغلقة تحت $f$ وكان $B
		\subseteq X$، فإن $\Closureofunder{f}{B} \subseteq X$.
	\end{enumerate}
\end{lem}

\textbf{تنبيه تصحيحي على الأصل (C054-01).} صُرّح بالمجموعة المحيطة $E$ وبكون
$f$ دالة من $E$ إلى نفسها؛ فبدون ذلك لا يلزم أن يكون تقاطع الجزئيات
منغلقًا تحت $f$، ولا يكون المتغير $X$ محصورًا في مجموعة معلومة.

\begin{proof}
يُسلك فيه بعينه مسلك \olref[sfr][infinite][dedekind]{closureproperties}.
\end{proof}

ولا يبقى قبل إثبات شرودر–برنشتاين إلا القضية الآتية:

\begin{prop}\ollabel{sbhelper}
إذا كان $A \subseteq B \subseteq C$ وكان $\cardeq{A}{C}$، فإن $\cardeq{A}{B}$ و$\cardeq{B}{C}$.
\end{prop}

\begin{proof}
إذا أعطيت !!a{bijection} $f \colon C \to A$، فبما أن $A\subseteq C$
ننظر إلى $f$ أيضًا بوصفها دالة من $C$ إلى نفسها، ولتكن $F =
\Closureofunder{f}{C \setminus B}$، وعرّف دالة $g$ مجالها $C$ كما
يأتي:
\[
	g(x) =
	\begin{cases}
			f(x) &\text{إذا كان $x \in F$}\\
			x & \text{فيما عدا ذلك}
		\end{cases}
\]
سنبيّن أن $g$ هي !!a{bijection} $C \to B$ (أي إن $\cardeq{B}{C}$)، ومن ذلك يلزم أن
$\comp{f^{-1}}{g} \colon A \to B$ هي !!a{bijection} (أي إن $\cardeq{A}{B}$)، فيكتمل البرهان.

نبدأ بتمييز الصورتين: إذا كان $x \in F$ وكان $y\notin F$، فلا بد أن $g(x) \neq
g(y)$. وإلا لزم $y = g(y) = g(x) =
f(x)$. لكن $x \in F$، و$F$ منغلقة تحت $f$ بمقتضى
\olref{Closureprops}؛ فيكون $y = f(x) \in F$، خلاف الفرض.

وعليه، إذا فرضنا $g(x) = g(y)$، لزم أن $x \in F$ إذا
وفقط إذا كان $y \in F$. فإن كان $x, y \in F$، كان $f(x) = g (x) =
g(y) = f(y)$، فيلزم $x = y$ لأن $f$ هي !!a{bijection}. وإن كان
$x, y \notin F$، حصل $x = g(x) = g(y) = y$. ففي الحالين يتساوى المدخلان؛ وبذلك تكون $g$
!!a{injection}.

ولإتمام البرهان نثبت $\ran{g} = B$. خذ $x \in B \subseteq C$. إن كان
$x \notin F$، فلدينا $g(x) = x$. وإن كان $x \in F$، فلا بد أن $x = f(y)$
لبعض $y \in F$. إذ لو انتفت هذه الصورة، لكانت $F \setminus \{x\}$
منغلقة تحت $f$ ومحتوية على $C\setminus B$، وهو محال بمقتضى
\olref{Closureprops}. وحينئذ $g(y) = f(y) = x$.
\end{proof}

فلنرجع إلى المبرهنة الأصلية. نذكّر بأن صورة الدالة $h$ على المجموعة $D$ معرّفة هكذا:
$\funimage{h}{D} = \Setabs{h(x)}{x
\in D}$.

\begin{proof}[برهان مبرهنة شرودر–برنشتاين] خذ !!{injection}s
$f \colon A \to B$ و$g \colon B \to A$. من $\funimage{f}{A}
\subseteq B$ يلزم $\funimage{g}{\funimage{f}{A}} \subseteq
g[B] \subseteq A$. والدالة $\comp{f}{g} \colon A \to
\funimage{g}{\funimage{f}{A}}$ هي !!{injection}، إذ كل من $g$ و
$f$ كذلك؛ ثم إن $\comp{f}{g}$ هي !!a{bijection} أيضًا، لاختيار صورتها مجالًا مقابلًا. فينتج
$\cardeq{\funimage{g}{\funimage{f}{A}}}{A}$، ومن
\olref{sbhelper} توجد !!a{bijection} $h \colon A \to
\funimage{g}{B}$. وأما $g^{-1}$ فهي !!a{bijection}
$\funimage{g}{B} \to B$. فتركيبهما $\comp{h}{g^{-1}} \colon A \to B$
!!a{bijection}، وهو المطلوب.
\end{proof}

\end{document}
